% Unit 096 terjemahan Bahasa Indonesia dari chapter7.tex baris 1661--1805.
% Sumber: Methods of Algebra, Volume 2, commit
% 9a5803ff77dd3257484cb177f851a73770a59dd3 (CC BY 4.0).
% stable-unit-id: o014.aljabr2.chapter7.morita-theory-a
% Cakupan: Bagian 7.7 dari pembukaan melalui Proposisi penurunan Morita;
% dilanjutkan oleh Unit 097 pada baris 1806.

% segment-id: o014.li2.chapter7.morita-theory-a.h001
\section{Teori Morita}\label{sec:Morita}
\hypertarget{o014.aljabr2.chapter7.morita-theory-a}{ }
% segment-id: o014.li2.chapter7.morita-theory-a.p001
Dalam bagian ini, tetapkan semesta Grothendieck $\mathcal{U}$ dan gelanggang
komutatif $\Bbbk$. Relatif terhadap $\mathcal{U}$, semua gelanggang secara
baku dianggap kecil. Funktor antara kategori $\Bbbk$-linear juga secara baku
dianggap $\Bbbk$-linear, sebagaimana dalam
Definisi~\ref{def:cat-k-linear} dan \ref{def:k-linear-cat}. Bila
$\Bbbk=\Z$, sifat $\Bbbk$-linear sama dengan sifat aditif.

% segment-id: o014.li2.chapter7.morita-theory-a.p002
Untuk kategori $\Bbbk$-linear $\mathcal{C}$ dan $\mathcal{D}$, notasi
$\mathrm{Fct}(\mathcal{C},\mathcal{D})$ atau
$\mathcal{D}^{\mathcal{C}}$ menyatakan kategori semua funktor
$F:\mathcal{C}\to\mathcal{D}$. Walaupun
$\mathrm{Fct}(\mathcal{C},\mathcal{D})$ belum tentu merupakan kategori-
$\mathcal{U}$, hal ini tidak menimbulkan kesulitan bagi pembahasan berikut
karena kita tidak membentuk limit atau konstruksi serupa di dalamnya.

% segment-id: o014.li2.chapter7.morita-theory-a.d001
\begin{definition}
	\index[sym1]{Fctc@$\mathrm{Fct}^c, \mathrm{Fct}_c$}
	Misalkan $\mathcal{C}$ dan $\mathcal{D}$ kategori $\Bbbk$-linear.
	\begin{itemize}
		\item Bila $\mathcal{C}$ dan $\mathcal{D}$ kokomplet, definisikan
		subkategori penuh $\mathrm{Fct}^c(\mathcal{C},\mathcal{D})$ dari
		$\mathrm{Fct}(\mathcal{C},\mathcal{D})$ yang terdiri atas semua funktor
		yang mempertahankan $\varinjlim$ kecil.
		\item Bila $\mathcal{C}$ dan $\mathcal{D}$ komplet, definisikan
		subkategori penuh $\cate{Fct}_c(\mathcal{C},\mathcal{D})$ dari
		$\mathrm{Fct}(\mathcal{C},\mathcal{D})$ yang terdiri atas semua funktor
		yang mempertahankan $\varprojlim$ kecil.
	\end{itemize}
	Dalam sebagian pustaka, objek $\mathrm{Fct}_c$ (atau
	$\mathrm{Fct}^c$) disebut funktor kontinu (atau kokontinu).
\end{definition}

% segment-id: o014.li2.chapter7.morita-theory-a.p003
Kita gunakan notasi dari \S\ref{sec:otimesL}. Untuk sembarang aljabar-$\Bbbk$
$A$ dan $B$, definisikan
% segment-id: o014.li2.chapter7.morita-theory-a.q001
\begin{align*}
	(A, B)\dcate{Mod} & := \text{kategori bimodul-$(A,B)$}, \\
	A\dcate{Mod} & := \text{kategori modul-$A$ kiri} \simeq (A, \Bbbk)\dcate{Mod}, \\
	\cated{Mod}B & := \text{kategori modul-$B$ kanan} \simeq (\Bbbk, B)\dcate{Mod}.
\end{align*}
Kategori-kategori tersebut merupakan contoh sederhana kategori
$\Bbbk$-linear.

% segment-id: o014.li2.chapter7.morita-theory-a.p004
Ingat kembali fakta dasar teori modul berikut. Setiap bimodul-$(A,B)$ $P$
secara kanonik menginduksi funktor-funktor yang mempertahankan
$\varinjlim$ kecil,
% segment-id: o014.li2.chapter7.morita-theory-a.q002
\[ P \dotimes{B} (\cdot): B\dcate{Mod} \to A\dcate{Mod}, \quad (\cdot) \dotimes{A} P : \cated{Mod}A \to \cated{Mod}B, \]
sedangkan setiap bimodul-$(B,A)$ $Q$ secara kanonik menginduksi funktor-
funktor yang mempertahankan $\varprojlim$ kecil,
% segment-id: o014.li2.chapter7.morita-theory-a.q003
\[ \Hom_{\cated{Mod}B}(Q, \cdot): B\dcate{Mod} \to A\dcate{Mod}, \quad \Hom_{\cated{Mod}A}(Q, \cdot): \cated{Mod}A \to \cated{Mod}B. \]
Jika sebagai gantinya diambil funktor $\Hom(\mathord\cdot,R)$ dengan $R$
bimodul-$(B,A)$, domainnya harus diganti dengan kategori lawan seperti
$(B\dcate{Mod})^{\opp}$ agar funktor itu tetap mempertahankan
$\varprojlim$ kecil.

% segment-id: o014.li2.chapter7.morita-theory-a.th001
\begin{theorem}[S.\ Eilenberg, C.\ E.\ Watts]\label{prop:Morita}
	Untuk sembarang aljabar-$\Bbbk$ $A$ dan $B$, terdapat ekuivalensi-
	ekuivalensi berikut:
	% segment-id: o014.li2.chapter7.morita-theory-a.q004
	\[\begin{tikzcd}[row sep=tiny]
		\mathrm{Fct}^c\left(B\dcate{Mod}, A\dcate{Mod}\right) & (A, B)\dcate{Mod} \arrow[l, "\sim"'] \arrow[r, "\sim"] & \mathrm{Fct}^c\left(\cated{Mod}A, \cated{Mod}B\right) \\
		P \dotimes{B} (\cdot) & P \arrow[mapsto, l] \arrow[mapsto, r] & (\cdot) \dotimes{A} P \\
		\mathrm{Fct}_c\left(B\dcate{Mod}, A\dcate{Mod}\right) & (B, A)\dcate{Mod} \arrow[l, "\sim"'] \arrow[r, "\sim"] & \mathrm{Fct}_c\left(\cated{Mod}A, \cated{Mod}B\right) \\
		\Hom_{B\dcate{Mod}}(Q, \cdot) & Q \arrow[mapsto, l] \arrow[mapsto, r] & \Hom_{\cated{Mod}A}(Q, \cdot) \\
		\mathrm{Fct}_c\left((B\dcate{Mod})^{\opp}, \cated{Mod}A \right) & (B, A)\dcate{Mod} \arrow[l, "\sim"'] \arrow[r, "\sim"] & \mathrm{Fct}_c\left((\cated{Mod}A)^{\opp}, B\dcate{Mod}\right) \\
		\Hom_{B\dcate{Mod}}(\cdot, R) & R \arrow[mapsto, l] \arrow[mapsto, r] & \Hom_{\cated{Mod}A}(\cdot, R).
	\end{tikzcd}\]

	Untuk ekuivalensi pertama dan kedua, komposisi funktor bersesuaian dengan
	hasil kali tensor bimodul; bila $A=B$, funktor identitas berasal dari $A$
	sebagai bimodul-$(A,A)$, hingga isomorfisme.
\end{theorem}
% segment-id: o014.li2.chapter7.morita-theory-a.pf001
\begin{proof}
	Argumen untuk berbagai ekuivalensi serupa. Kita hanya membahas
	% segment-id: o014.li2.chapter7.morita-theory-a.q005
	\[ (A, B)\dcate{Mod} \to \mathrm{Fct}^c(B\dcate{Mod}, A\dcate{Mod}). \]
	Pembahasan sebelum teorema telah menunjukkan bahwa
	$P\dotimes{B}(\mathord\cdot)$ memang objek dari
	$\mathrm{Fct}^c(B\dcate{Mod},A\dcate{Mod})$, dan homomorfisme bimodul
	$P\to P'$ jelas menginduksi morfisme funktor
	$P\dotimes{B}(\mathord\cdot)\to P'\dotimes{B}(\mathord\cdot)$. Sekarang
	kita konstruksikan funktor kuasi-inversnya.

	Untuk objek $F$ dari
	$\mathrm{Fct}^c(B\dcate{Mod},A\dcate{Mod})$, pandang $B$ sebagai modul-$B$
	kiri dan definisikan $P:=F(B)$. Karena $B$ juga beraksi pada dirinya dari
	kanan melalui perkalian dan $F$ adalah funktor $\Bbbk$-linear, $P$ secara
	alami memperoleh struktur bimodul-$(A,B)$. Isomorfisme kanonik
	$P\dotimes{B}B\simeq P$ menunjukkan bahwa
	% segment-id: o014.li2.chapter7.morita-theory-a.q006
	\[ (A, B)\dcate{Mod} \to \mathrm{Fct}^c\left(B\dcate{Mod}, A\dcate{Mod} \right) \xrightarrow{F \mapsto P} (A, B)\dcate{Mod} \;\text{merupakan komposisi}\; \simeq \identity. \]

	Untuk komposisi dalam arah sebaliknya, ambil funktor $F$ seperti di atas
	dan tetapkan $P:=F(B)$. Untuk sembarang modul-$B$ kiri $M$, terdapat
	homomorfisme surjektif kanonik
	% segment-id: o014.li2.chapter7.morita-theory-a.q007
	\[ \bigoplus_{\varphi \in \Hom(B, M)} B \twoheadrightarrow M. \]
	Perhatikan bahwa $\bigoplus_\varphi$ bersifat ``kecil''. Dengan mengulangi
	konstruksi tersebut pada kernel homomorfisme surjektif itu, diperoleh
	barisan eksak kanonik untuk $M$,
	% segment-id: o014.li2.chapter7.morita-theory-a.q008
	\begin{equation}\label{eqn:Morita-aux0}
		\bigoplus_\psi B \to \bigoplus_\varphi B \to M \to 0,
	\end{equation}
	dan bagian pertamanya dapat dipandang sebagai perkalian kanan dengan
	matriks tak hingga yang bernilai dalam
	$B=\End_{B\dcate{Mod}}(B)$. Karena $F$ dan
	$P\dotimes{B}(\mathord\cdot)$ keduanya mempertahankan $\varinjlim$ kecil,
	menerapkan keduanya pada \eqref{eqn:Morita-aux0} memberikan barisan eksak
	dalam $A\dcate{Mod}$,
	% segment-id: o014.li2.chapter7.morita-theory-a.q009
	\begin{gather*}
		\bigoplus_\psi P \to \bigoplus_\varphi P \to F(M) \to 0, \\
		\bigoplus_\psi P \to \bigoplus_\varphi P \to P \dotimes{B} M \to 0.
	\end{gather*}
	Dalam kedua baris, pemetaan
	$\bigoplus_\psi P\to\bigoplus_\varphi P$ berasal dari matriks yang sama
	(dengan entri dalam $B$), sehingga diperoleh isomorfisme kanonik
	$F(M)\simeq P\dotimes{B}M$.

Berdasarkan kendala asosiativitas hasil kali tensor, ekuivalensi pertama
	$(A,B)\dcate{Mod}\to\mathrm{Fct}^c(B\dcate{Mod},A\dcate{Mod})$ membawa
hasil kali tensor bimodul ke komposisi funktor, hingga isomorfisme. Untuk
	ekuivalensi kedua, pernyataan terkait mengikuti dari adjungsi antara produk
	tensor dan $\Hom$; lihat \cite[Teorema 6.6.5]{Li1}. Dalam kedua kasus,
	mudah dilihat bahwa bimodul-$(A,A)$ $A$ bersesuaian dengan funktor identitas,
	hingga isomorfisme.
\end{proof}

% segment-id: o014.li2.chapter7.morita-theory-a.p005
Sekarang ukuran kategori-kategori funktor tersebut mudah dikendalikan.

% segment-id: o014.li2.chapter7.morita-theory-a.c001
\begin{corollary}
	Untuk sembarang aljabar-$\Bbbk$ $A$ dan $B$, semua kategori funktor
	$\mathrm{Fct}^c(\cdots)$ dan $\mathrm{Fct}_c(\cdots)$ yang muncul dalam
	Teorema~\ref{prop:Morita} adalah kategori-$\mathcal{U}$.
\end{corollary}
% segment-id: o014.li2.chapter7.morita-theory-a.pf002
\begin{proof}
	Sebab $(A,B)\dcate{Mod}$ dan $(B,A)\dcate{Mod}$ keduanya kategori-
	$\mathcal{U}$.
\end{proof}

% segment-id: o014.li2.chapter7.morita-theory-a.dp001
\begin{definition-proposition}\label{def:Morita-equiv}
	\index{ekuivalensi Morita@ekuivalensi Morita (Morita equivalence)}
	Untuk sembarang aljabar-$\Bbbk$ $A$ dan $B$, pernyataan-pernyataan berikut
	ekuivalen.
	\begin{enumerate}[(i)]
		\item $A\dcate{Mod}$ dan $B\dcate{Mod}$ ekuivalen.
		\item Terdapat bimodul-$(A,B)$ $P$ dan bimodul-$(B,A)$ $Q$ serta
		\begin{compactitem}
			\item isomorfisme bimodul-$(A,A)$
			$P\dotimes{B}Q\simeq A$;
			\item isomorfisme bimodul-$(B,B)$
			$Q\dotimes{A}P\simeq B$.
		\end{compactitem}
		\item $\cated{Mod}A$ dan $\cated{Mod}B$ ekuivalen.
	\end{enumerate}
	Bila salah satu kondisi di atas berlaku, $A$ dan $B$ disebut
	\emph{ekuivalen Morita}.
\end{definition-proposition}
% segment-id: o014.li2.chapter7.morita-theory-a.pf003
\begin{proof}
	Ekuivalensi niscaya mempertahankan $\varinjlim$ dan $\varprojlim$.
	Teorema~\ref{prop:Morita} memberikan (i) $\iff$ (ii) dan
	(iii) $\iff$ (ii).
\end{proof}

% segment-id: o014.li2.chapter7.morita-theory-a.p006
Kondisi (ii) ringkas tetapi belum cukup eksplisit. Kondisi itu akan
diperhalus oleh Teorema~\ref{prop:Morita-refined} pada akhir bagian ini.
Sebelumnya, beberapa contoh dan teori terkait perlu diperkenalkan.

% segment-id: o014.li2.chapter7.morita-theory-a.eg001
\begin{example}
	Untuk aljabar-$\Bbbk$ komutatif, ekuivalensi Morita sama artinya dengan
	isomorfisme aljabar. Hal ini didasarkan pada isomorfisme kanonik
	$Z(A\dcate{Mod})\simeq Z(A)$, dengan ruas kiri pusat kategori abelian dan ruas
	kanan pusat aljabar-$\Bbbk$.
\end{example}

% segment-id: o014.li2.chapter7.morita-theory-a.eg002
\begin{example}
	Misalkan $n\in\Z_{\geq1}$. Setiap aljabar-$\Bbbk$ $A$ ekuivalen Morita
	dengan aljabar matriks $n\times n$, $\mathrm{M}_n(A)$. Untuk melihat hal
	ini, definisikan melalui operasi matriks
	% segment-id: o014.li2.chapter7.morita-theory-a.q010
	\begin{align*}
		P & := \text{bimodul-$(A,\mathrm{M}_{n \times n}(A))$ } A^n \quad \text{(vektor baris)}, \\
		Q & := \text{bimodul-$(\mathrm{M}_{n \times n}(A),A)$ } A^n \quad \text{(vektor kolom)}.
	\end{align*}
	Perkalian matriks memberikan isomorfisme yang diperlukan,
	% segment-id: o014.li2.chapter7.morita-theory-a.q011
	\begin{gather*}
		P \dotimes{\mathrm{M}_{n \times n}(A)} Q \rightiso A \quad \text{(sebagai bimodul-$(A,A)$)}, \\
		Q \dotimes{A} P \rightiso \mathrm{M}_{n \times n}(A) \quad \text{(sebagai bimodul-$(\mathrm{M}_{n \times n}(A), \mathrm{M}_{n \times n}(A))$)}.
	\end{gather*}
\end{example}

% segment-id: o014.li2.chapter7.morita-theory-a.p007
Kembali ke keadaan umum. Untuk sembarang bimodul-$(A,B)$ $P$, fakta dasar
teori modul \cite[Teorema 6.6.5]{Li1} memberikan pasangan adjoin
% segment-id: o014.li2.chapter7.morita-theory-a.q012
\begin{equation}\label{eqn:AB-adjunction}\begin{tikzcd}
	(\cdot) \dotimes{A} P : \cated{Mod}A \arrow[shift left, r] & \cated{Mod}B: \Hom_{\cated{Mod}B}(P, \cdot). \arrow[shift left, l]
\end{tikzcd}\end{equation}
Teorema~\ref{prop:Morita} menunjukkan bahwa, hingga isomorfisme, funktor di
ruas kiri dan kanan masing-masing mencakup semua objek
$\mathrm{Fct}^c(\cated{Mod}A,\cated{Mod}B)$ dan
$\mathrm{Fct}_c(\cated{Mod}B,\cated{Mod}A)$.

% segment-id: o014.li2.chapter7.morita-theory-a.p008
Berdasarkan Contoh~\ref{eg:adjunction-monad}, pasangan adjoin
\eqref{eqn:AB-adjunction} menentukan monad $T$ pada $\cated{Mod}A$ dan
komonad $L$ pada $\cated{Mod}B$. Deskripsi umumnya sebagai berikut. Tinjau
modul-$A$ kanan $N$ dan modul-$B$ kanan $M$:
% segment-id: o014.li2.chapter7.morita-theory-a.q013
\begin{equation}\label{eqn:bimod-adjunction}
	\begin{array}{|c|c|} \hline
		\text{unit } \eta_N & \text{kounit } \epsilon_M \\ \hline
		N \to \Hom_{\cated{Mod}B}\left( P, N \dotimes{A} P \right) & \Hom_{\cated{Mod}B}(P, M) \dotimes{A} P \to M \\
		x \mapsto \left[ p \mapsto x \otimes p \right] & \varphi \otimes p \mapsto \varphi(p) \\ \hline
	\end{array}
\end{equation}

% segment-id: o014.li2.chapter7.morita-theory-a.p009
Ingat kembali definisi funktor eksak dalam Definisi~\ref{def:exact-functor}.

% segment-id: o014.li2.chapter7.morita-theory-a.pr001
\begin{proposition}\label{prop:Morita-descent}
	Misalkan $P$ bimodul-$(A,B)$. Jika
	$(\mathord\cdot)\dotimes{A}P$ (atau
	$\Hom_{\cated{Mod}B}(P,\mathord\cdot)$) merupakan funktor eksak setia,
	maka pasangan adjoin \eqref{eqn:AB-adjunction} bersifat komonadik (atau
	monadik); lihat Definisi~\ref{def:monadic} dan versi dualnya.
\end{proposition}
% segment-id: o014.li2.chapter7.morita-theory-a.pf004
\begin{proof}
	Berikut bukti untuk kasus $(\mathord\cdot)\dotimes{A}P$; argumen untuk
	sisi lain persis sama.

	Telah ditunjukkan bahwa $(\mathord\cdot)\dotimes{A}P$ mempunyai adjoin
	kanan. Tinggal memeriksa versi dual syarat dalam Teorema~\ref{prop:Beck}
	(iii). Pertama, kita buktikan bahwa $(\mathord\cdot)\dotimes{A}P$ adalah
	funktor konservatif. Untuk homomorfisme modul-$A$ kanan $f:N\to N'$,
	konservativitas yang diperlukan mengikuti dari
	% segment-id: o014.li2.chapter7.morita-theory-a.q014
	\begin{multline*}
		f \;\text{adalah isomorfisme} \iff 0 \to N \xrightarrow{f} N' \to 0 \;\text{eksak} \\
		\stackrel{\text{Proposisi \ref{prop:faithful-exact} (iii)}}{\iff} 0 \to N \dotimes{A} P \xrightarrow{\identity \otimes f} N' \dotimes{A} P \to 0 \;\text{eksak} \iff \identity \otimes f \;\text{adalah isomorfisme}.
	\end{multline*}

	Selanjutnya, setiap pasangan morfisme dalam $\cated{Mod}A$ dan
	$\cated{Mod}B$ mempunyai penyama, dan hipotesis eksak setia memastikan
	bahwa $(\mathord\cdot)\dotimes{A}P$ mempertahankan penyama. Jadi, seluruh
	syarat Teorema~\ref{prop:Beck} (iii) terpenuhi.
\end{proof}

% segment-id: o014.li2.chapter7.morita-theory-a.p010
Ini baru suatu hasil abstrak. Untuk menerapkannya, monad $T$ dan komonad $L$
yang bersesuaian perlu dinyatakan secara eksplisit dalam beberapa kasus
khusus. Kita mulai dengan perubahan gelanggang.
